\chapter{Was NOMAD nicht kann}

Damit Sie sich im Notfall nicht auf etwas verlassen, das nicht da ist, hier die bekannten Grenzen dieser deutschen Fassung (Stand Oktober 2026).

\begin{description}[leftmargin=1.2em,style=nextline]
\item[Manche Inhalte sind nur auf Englisch da.] Oberfläche, Installationsprogramm, Anleitungen und dieses Heft sind deutsch. Es gibt deutsche Wikipedia, Medizin-Enzyklopädie (WikiMed), Lehrbücher, Wörterbuch, Koch- und Reparaturanleitungen. Für \emph{Überleben und Vorsorge} und die medizinischen Nachschlagewerke (CDC, NHS, Militärmedizin) gibt es bisher nur englische Bücher. Deshalb sind die englischen Kategorien mit dem Zusatz „Englisch“ gekennzeichnet.
\item[Karten: Ortsnamen und Länderliste englisch.] Karten für Deutschland und jedes andere Land sind möglich (Kapitel~5). Die Auswahlliste der Länder und die Beschriftung der Karte („Germany“, „Cologne“) sind aber englisch. Die „kuratierten Kartenregionen“ im Karten-Manager enthalten nur Gebiete der USA; für Deutschland nehmen Sie „Länder auswählen“.
\item[Kein Windows-Programm.] NOMAD läuft unter Linux. Auf Windows ohne Linux gibt es nur eine Anleitung des Originals (WSL2) in englischer Sprache. Dieses Heft empfiehlt deshalb die USB-SSD.
\item[Kein Mac, keine ARM-Rechner.] Das System ist für gängige PCs (x86\_64) gebaut. Macs mit Apple-Chip und kleine ARM-Platinen (zum Beispiel Raspberry Pi) werden von dieser Version nicht unterstützt.
\item[Kein Passwortzwang.] Das System der USB-SSD hat bei jeder SSD dasselbe Anfangspasswort. Ändern Sie es (Kapitel~2).
\item[Fremde Hardware.] Ob das System auf Ihrem PC startet (Grafik, WLAN, Startmenü), hängt vom PC ab und lässt sich nicht für alle Geräte vorab zusichern.
\item[Updates.] NOMAD kann sich selbst aktualisieren, aber nur mit Internet; Inhalte veralten (Datum in Kapitel~9 notieren).
\end{description}
